% ============================================================================
% Fichier  : partie3/P03_C12_ethique.tex
% Titre    : Éthique et Responsabilité
% Auteur   : MINKA MI NGUIDJOÏ Thierry Emmanuel
% Version  : 1.0 -- novembre 2026
% ============================================================================
% Sources : code d'éthique de l'ISC2 (quatre canons) ; code de déontologie
% professionnelle de l'ISACA ; code d'éthique de l'ACM (2018) ; ISO/IEC
% 29147:2018 et ISO/IEC 30111:2019 (divulgation et traitement des
% vulnérabilités) ; loi n° 2010/012, art. 44 et 65 (texte officiel).
% ============================================================================

\chapter{Éthique et Responsabilité}
\label{chap:ethique}
\chaptermark{Éthique}

\epigraph{%
    « Celui qui détient les clés n'est pas propriétaire de la maison. »%
}{Principe du cours}

\niveauChapitre{Fondamental}{Chapitres~\ref{chap:gouvernance}
et~\ref{chap:droit}.}

\begin{boxobjectifs}
\begin{itemize}
    \item Expliquer pourquoi les métiers de la sécurité ont besoin d'une
          éthique propre.
    \item Présenter les codes de l'ISC2, de l'ISACA et de l'ACM.
    \item Conduire une divulgation coordonnée de vulnérabilité.
    \item Résoudre un dilemme avec une méthode écrite.
    \item Situer les limites légales et éthiques des tests et de la
          surveillance.
\end{itemize}
\end{boxobjectifs}

\bigskip

% ============================================================================
\section*{Question directrice : que doit faire celui qui a tous les droits~?}
% ============================================================================

L'administrateur de LiZbiZ peut lire chaque courriel, chaque fichier,
chaque journal. Pendant l'enquête sur l'Opération GAROUA, le directeur
lui demande de « regarder les messages » d'un salarié qu'il soupçonne. Il
en a la capacité technique~; en a-t-il le droit, et doit-il le faire~?
Les métiers de la sécurité donnent un pouvoir réel sur la vie des autres~:
ils exigent une éthique explicite.

\medskip
\textbf{Plan du chapitre.}
\begin{enumerate}
    \item Pourquoi une éthique de la sécurité (section~\ref{sec:c12-pourquoi}).
    \item Les codes professionnels (section~\ref{sec:c12-codes}).
    \item La divulgation des vulnérabilités (section~\ref{sec:c12-divulgation}).
    \item Résoudre un dilemme (section~\ref{sec:c12-dilemme}).
    \item Les limites des tests et de la surveillance (section~\ref{sec:c12-limites}).
\end{enumerate}

\separateur

% ============================================================================
\section{Pourquoi une Éthique de la Sécurité}
\label{sec:c12-pourquoi}
% ============================================================================

Les professionnels de la sécurité ont des accès privilégiés, des
compétences offensives et des informations sur les faiblesses des
autres. Les mêmes gestes servent à protéger ou à nuire~: un balayage de
ports, la lecture d'un journal, le contournement d'une protection. Ce qui
les distingue est le mandat, le but et la manière. C'est pourquoi les
organisations professionnelles ont adopté des codes, dont la violation peut
entraîner le retrait d'une certification.

% ============================================================================
\section{Les Codes Professionnels}
\label{sec:c12-codes}
% ============================================================================

\begin{tableaucours}{L{2.4cm}L{9.8cm}}{%
    \tableauHeader{Code} &
    \tableauHeader{Principes essentiels}
}
ISC2 & Quatre canons~: protéger la société, le bien commun et
l'infrastructure~; agir honorablement, honnêtement, justement,
responsablement et légalement~; rendre un service diligent et compétent
aux mandants~; faire progresser et protéger la profession~\cite{isc2ethics} \\
ISACA & Soutenir les normes~; objectivité, diligence~; légalité~;
confidentialité~; compétence~; informer les parties
concernées~\cite{isaca2024ethique} \\
ACM & Contribuer au bien de la société~; éviter de nuire~; honnêteté~;
respect de la vie privée~; confidentialité~; qualité du
travail~\cite{acm2018code} \\
\end{tableaucours}
\vspace{-0.8em}
\begin{center}{\small\itshape Tableau~12.1. Trois codes d'éthique}\end{center}

Les trois codes convergent~: légalité, honnêteté, compétence,
confidentialité, et une responsabilité qui dépasse le seul employeur
(la société, les personnes concernées).

% ============================================================================
\section{La Divulgation des Vulnérabilités}
\label{sec:c12-divulgation}
% ============================================================================

\begin{boxdef}
\termecle{Divulgation coordonnée}~: processus par lequel la personne qui
découvre une vulnérabilité la signale d'abord à l'éditeur ou à
l'exploitant, lui laisse un délai raisonnable pour corriger, puis la rend
publique de façon concertée. L'ISO/IEC~29147 décrit la réception des
signalements par un éditeur~\cite{iso29147_2018}~; l'ISO/IEC~30111 leur
traitement interne~\cite{iso30111_2019}.
\end{boxdef}

\begin{tableaucours}{L{3.4cm}L{8.8cm}}{%
    \tableauHeader{Étape} &
    \tableauHeader{Bonne pratique}
}
Découverte & Ne pas aller au-delà de ce qui prouve la vulnérabilité~;
ne pas consulter ni copier de données \\
Signalement & À l'éditeur ou à l'exploitant, par le canal prévu (adresse
dédiée, CERT national) \\
Délai & Convenir d'un délai de correction (souvent de l'ordre de 90~jours) \\
Correction & L'éditeur corrige et informe ses clients \\
Publication & Description concertée, sans code d'exploitation prêt à
l'emploi avant la correction \\
\end{tableaucours}
\vspace{-0.8em}
\begin{center}{\small\itshape Tableau~12.2. Les étapes de la divulgation coordonnée}\end{center}

\begin{boxalert}
\textbf{Violation typique.} Démontrer une faille du portail de l'ANPP en
téléchargeant la liste des bénéficiaires, « pour être pris au sérieux ».

\smallskip
\textbf{Règle~:} la preuve s'arrête au minimum nécessaire. Consulter ou
copier des données, c'est passer de la découverte à l'atteinte, avec ses
conséquences pénales (Chapitre~\ref{chap:droit}).
\end{boxalert}

% ============================================================================
\section{Résoudre un Dilemme}
\label{sec:c12-dilemme}
% ============================================================================

\begin{boxprocedure}
\textbf{Une méthode en six étapes}
\begin{enumerate}
    \item \textbf{Établir les faits}, séparés des suppositions.
    \item \textbf{Identifier les parties} concernées et leurs intérêts.
    \item \textbf{Repérer les règles}~: loi, politique interne, code
          professionnel, contrat.
    \item \textbf{Lister les options}, y compris refuser ou demander un
          avis.
    \item \textbf{Évaluer chaque option} au regard des règles et des
          conséquences pour chaque partie.
    \item \textbf{Décider et documenter}, par écrit, la décision et ses
          raisons.
\end{enumerate}
\end{boxprocedure}

\begin{boxscenario}{Opération GAROUA. « Regardez ses messages »}
\textbf{Faits.} Le directeur soupçonne un comptable d'avoir aidé
l'attaquant, sans élément précis. Il demande à l'administrateur de lire
sa messagerie.

\medskip
\textbf{Règles.} La loi interdit d'intercepter ou de stocker des
communications sans consentement, sauf autorisation légale (art.~44,
Chapitre~\ref{chap:droit}). La politique de LiZbiZ ne prévoit aucune
procédure d'accès à une boîte personnelle.

\medskip
\textbf{Décision} (confiance élevée). L'administrateur refuse de lire les
messages de sa propre initiative, conserve les journaux de connexion qui
relèvent de l'enquête, et renvoie le directeur vers le conseil juridique
et, s'il y a lieu, vers une plainte qui permettra aux enquêteurs d'agir
dans un cadre légal. Il consigne la demande et sa réponse.
\end{boxscenario}

% ============================================================================
\section{Les Limites des Tests et de la Surveillance}
\label{sec:c12-limites}
% ============================================================================

\begin{tableaucours}{L{4.0cm}L{8.2cm}}{%
    \tableauHeader{Activité} &
    \tableauHeader{Condition éthique et légale}
}
Test d'intrusion & Autorisation écrite du responsable du système~;
périmètre et fenêtre définis (Chapitre~\ref{chap:evaluer}) \\
Hameçonnage simulé & Approbation de la direction~; pas de sanction
individuelle (Chapitre~\ref{chap:humain}) \\
Surveillance des journaux & Finalité de sécurité déclarée~; accès limité~;
durée de conservation fixée \\
Accès à une boîte de messagerie & Cadre légal et procédure écrite~; jamais
à l'initiative de l'administrateur \\
Recherche de vulnérabilités sur Internet & Seulement sur ses propres
systèmes ou avec autorisation \\
\end{tableaucours}
\vspace{-0.8em}
\begin{center}{\small\itshape Tableau~12.3. Les conditions des activités sensibles}\end{center}

\begin{boxmetier}{Administrateur systèmes et réseaux}
\textbf{Le pouvoir de l'administrateur.} L'administrateur signe une charte
qui rappelle ses obligations de confidentialité et ce qu'il ne doit pas
faire sans procédure. Il trace ses propres accès aux données sensibles.
\end{boxmetier}

\begin{boxmetier}{Auditeur de sécurité}
\textbf{Rester dans le mandat.} L'auditeur ne teste que ce que sa lettre
de mission autorise, ne conserve que les preuves nécessaires, et ne
divulgue ses constats qu'au commanditaire.
\end{boxmetier}

\separateur

\begin{boxresume}
\textbf{Ce qu'il faut retenir de ce chapitre~:}
\begin{itemize}
    \item \textbf{Les métiers de la sécurité} donnent un pouvoir qui exige
          une éthique explicite~: le mandat, le but et la manière font la
          différence.
    \item \textbf{Les codes} de l'ISC2, de l'ISACA et de l'ACM convergent~:
          légalité, honnêteté, compétence, confidentialité, responsabilité
          envers la société.
    \item \textbf{La divulgation coordonnée} (ISO/IEC~29147 et~30111)
          protège les utilisateurs et le découvreur.
    \item \textbf{Un dilemme} se résout par une méthode écrite et
          documentée.
    \item \textbf{Tests et surveillance} exigent autorisation, finalité et
          procédure.
\end{itemize}
\end{boxresume}

\begin{boxinfo}
\textbf{Points de contrôle de l'élève.} Je sais~: citer les quatre canons
de l'ISC2~; décrire la divulgation coordonnée~; appliquer la méthode en
six étapes~; dire à quelles conditions un test ou une surveillance est
admissible.
\end{boxinfo}

\bigskip

\begin{boxexo}[title={\ding{45}\quad Exercice~12.1. Le même geste \niveaubase}]
Pour trois gestes techniques (balayage de ports, lecture d'un journal,
contournement d'un mot de passe), décrivez un usage légitime et un usage
illégitime.
\end{boxexo}

\begin{boxexo}[title={\ding{45}\quad Exercice~12.2. Les canons \niveaubase}]
Rattachez à un canon de l'ISC2~: refuser une mission qu'on ne sait pas
conduire~; signaler une faille qui menace un hôpital~; ne pas mentir sur
ses certifications.
\end{boxexo}

\begin{boxexo}[title={\ding{45}\quad Exercice~12.3. La découverte \niveaumoyen}]
Un étudiant découvre que le portail de résultats de son école affiche les
notes de tous les étudiants en changeant un numéro dans l'adresse. Que
doit-il faire, et ne pas faire~?
\end{boxexo}

\begin{boxexo}[title={\ding{45}\quad Exercice~12.4. Le délai \niveaumoyen}]
L'éditeur ne répond pas après 90~jours. Quelles options restent au
découvreur, et lesquelles sont responsables~?
\end{boxexo}

\begin{boxexo}[title={\ding{45}\quad Exercice~12.5. La méthode \niveaumoyen}]
Appliquez la méthode en six étapes à ce cas~: un client important de
LiZbiZ demande à l'administrateur la liste des expéditions d'un
concurrent, « contre un geste ».
\end{boxexo}

\begin{boxexo}[title={\ding{45}\quad Exercice~12.6. Surveiller les salariés \niveaumoyen}]
La BANQUE SANAGA veut enregistrer tous les sites consultés par ses
agents. À quelles conditions est-ce acceptable~?
\end{boxexo}

\begin{boxexo}[title={\ding{45}\quad Exercice~12.7. Le conflit d'intérêts \niveaumoyen}]
Le RSSI de LiZbiZ est aussi actionnaire de SAHEL INFOGÉRANCE. Analysez la
situation et proposez une solution.
\end{boxexo}

\begin{boxexo}[title={\ding{45}\quad Exercice~12.8. La charte \niveaumoyen}]
Rédigez la charte d'un administrateur systèmes~: obligations, interdits,
traçabilité de ses accès.
\end{boxexo}

\begin{boxexo}[title={\ding{45}\quad Exercice~12.9. Une politique de divulgation \niveauexpert}]
Rédigez la politique de divulgation de vulnérabilités de LiZbiZ, publiée
sur son site~: canal, engagements, délais, ce que LiZbiZ demande aux
découvreurs.
\end{boxexo}

\begin{boxexo}[title={\ding{45}\quad Exercice~12.10. Le dilemme du RSSI \niveauexpert}]
Le directeur général demande au RSSI de ne pas signaler l'incident au
conseil d'administration avant la clôture des comptes. Appliquez la
méthode et rédigez la réponse du RSSI.
\end{boxexo}

\bigskip

\begin{boxinfo}
\textbf{Pour aller plus loin.}
\begin{itemize}
    \item Les codes de l'ISC2, de l'ISACA et de
          l'ACM~\cite{isc2ethics,isaca2024ethique,acm2018code}.
    \item ISO/IEC~29147 et 30111~\cite{iso29147_2018,iso30111_2019}.
\end{itemize}
\end{boxinfo}

\begin{boxinfo}
\textbf{Transition.} La Partie~\ref{part:technique} passe à la sécurité
technique, en commençant par les modèles formels qui fondent le contrôle
d'accès.
\end{boxinfo}
